\documentclass[aspectratio=169,11pt]{beamer}

% Theme and color setup
\usetheme{default}
\usecolortheme{default}

% Remove navigation symbols
\setbeamertemplate{navigation symbols}{}

% Custom colors (MIT-inspired)
\definecolor{mitred}{RGB}{163,31,52}
\definecolor{darkblue}{RGB}{0,40,85}
\definecolor{lightgray}{RGB}{245,245,245}
\definecolor{darkgray}{RGB}{64,64,64}

% Set main colors
\setbeamercolor{frametitle}{fg=white,bg=darkblue}
\setbeamercolor{title}{fg=darkblue}
\setbeamercolor{author}{fg=black}
\setbeamercolor{institute}{fg=black}
\setbeamercolor{date}{fg=black}
\setbeamercolor{structure}{fg=darkblue}
\setbeamercolor{normal text}{fg=black}
\setbeamercolor{itemize item}{fg=mitred}
\setbeamercolor{itemize subitem}{fg=darkblue}
\setbeamercolor{enumerate item}{fg=darkblue}

% Frame title band
\newlength{\titlebandht}
\setlength{\titlebandht}{4ex}

\setbeamertemplate{frametitle}{
  \nointerlineskip
  \begin{beamercolorbox}[wd=\paperwidth,ht=\titlebandht,dp=1.1ex,leftskip=1.2em,rightskip=1.2em]{frametitle}
    \usebeamerfont{frametitle}\insertframetitle
    \ifx\insertframesubtitle\@empty\relax\else\\[-0.2ex]
      \usebeamerfont{framesubtitle}\insertframesubtitle
    \fi
  \end{beamercolorbox}
  \vspace{0.6em}
}

% Font settings
\usefonttheme{professionalfonts}
\usefonttheme{serif}

% Footline with page numbers
\setbeamertemplate{footline}{
    \hfill\insertframenumber\hspace{2em}\vspace{0.5em}
}

% Packages
\usepackage{amsmath,amssymb,amsthm}
\usepackage{graphicx}
\usepackage{tikz}
\usepackage{booktabs}
\usepackage{mathtools}
\usepackage{bm}
\usepackage{bbm}
\usepackage{hyperref}

% Custom commands for math
\newcommand{\E}{\mathbb{E}}
\renewcommand{\P}{\mathbb{P}}
\newcommand{\R}{\mathbb{R}}
\newcommand{\N}{\mathbb{N}}
\newcommand{\Z}{\mathbb{Z}}
\newcommand{\F}{\mathcal{F}}
\newcommand{\Cov}{\text{Cov}}
\newcommand{\Var}{\text{Var}}
\newcommand{\Corr}{\text{Corr}}
\newcommand{\rank}{\mathrm{rank}}
\DeclareMathOperator*{\argmin}{arg\,min}
\DeclareMathOperator*{\argmax}{arg\,max}
\DeclareMathOperator*{\plim}{plim}

% Block colors
\setbeamercolor{block title}{bg=lightgray,fg=darkblue}
\setbeamercolor{block body}{bg=lightgray!50}

% Alert block colors
\setbeamercolor{block title alerted}{bg=mitred!15,fg=mitred}
\setbeamercolor{block body alerted}{bg=mitred!5}

% Title information
\title{Problem Set 4: Solutions}
\subtitle{SVARs: Short-Run and Long-Run Identification, and VECMs}
\author{Rafael Lincoln}
\institute{Econometrics II -- NYU PhD Program}
\date{Spring 2026}

\begin{document}

% ============================================================================
% TITLE SLIDE
% ============================================================================
\begin{frame}[plain]
    \titlepage
\end{frame}

% ============================================================================
% OUTLINE
% ============================================================================
\begin{frame}{Today's Roadmap}
    \textbf{Problem Set 4 Solutions}
    \begin{enumerate}
        \item \textbf{Introduction: The Identification Problem in Macroeconometrics}
        \begin{itemize}
            \item SVMA representation, fundamental non-identification, rotation parameterization
        \end{itemize}
        \vspace{0.3em}
        \item \textbf{Problem 1: Gal\'{\i} (1999) --- Long-Run and Short-Run Restrictions}
        \begin{itemize}
            \item Long-run productivity restriction, CEE, identification analysis, computational procedure
        \end{itemize}
        \vspace{0.3em}
        \item \textbf{Problem 2: SVECM --- Cointegration and Common Trends}
        \begin{itemize}
            \item Granger representation, structural long-run impact, noninvertibility of $\Delta x_t$, SVECM vs VAR in levels
        \end{itemize}
        \vspace{0.3em}
        \item \textbf{Problem 3: CFP Critique --- Cholesky Under Misspecification}
        \begin{itemize}
            \item Identified shocks as mixtures, sufficient condition for sign reversal, robustness
        \end{itemize}
    \end{enumerate}
\end{frame}

% ============================================================================
% INTRODUCTION: THE IDENTIFICATION PROBLEM
% ============================================================================
\begin{frame}[plain]
    \begin{center}
        {\Large \textbf{Introduction}}\\[1em]
        {\large The Identification Problem in Macroeconometrics}
    \end{center}
\end{frame}

% ----------------------------------------------------------------------------
\begin{frame}{Motivation}
    \textbf{Objects of interest.} Most Macroeconomic Models can be expressed under SVMA representation:
    \[
    y_t = \sum_{\ell=0}^{\infty} \Theta_\ell\, \varepsilon_{t-\ell}, \qquad \varepsilon_t \sim WN(0, I_n).
    \]
    \begin{itemize}
        \item From $\{\Theta_\ell\}$ we get: \textbf{IRFs}, \textbf{FVDs}, \textbf{Historical Decompositions}
        \item These are the key objects for policy analysis and testing macro models
    \end{itemize}

    \vspace{0.5em}

    \textbf{How can we learn $\{\Theta_\ell\}$ from data?} Three broad approaches:
    \begin{enumerate}
        \item \textbf{Semi-structural, time series}: use time-series properties of $y_t$ + identifying assumptions for \textit{families} of models \textcolor{mitred}{$\leftarrow$ our focus today}
        \item \textbf{Semi-structural, cross-sectional}: microeconometric causal designs
        \item \textbf{Full structural}: specify and estimate a DSGE model
    \end{enumerate}
\end{frame}

% ----------------------------------------------------------------------------
\begin{frame}{The Fundamental Identification Problem}
    \begin{block}{Proposition (Non-identification from second moments)}
        Under standard assumptions, we can find $\{\Theta_\ell^*\}$ and $\{\varepsilon_t^*\}$ such that $\Theta_\ell \neq \Theta_\ell^*$ for some $\ell$, $\varepsilon_t^* \sim WN(0, I_n)$, and for all $h$:
        \[
        \Gamma_h = \Cov\!\left(\sum_{\ell=0}^\infty \Theta_\ell \varepsilon_{t-\ell},\; \sum_{\ell=0}^\infty \Theta_\ell \varepsilon_{t-\ell-h}\right) = \sum_{\ell=0}^\infty \Theta_\ell^*\,(\Theta_{\ell+h}^*)'.
        \]
    \end{block}

    \vspace{0.3em}
    \textbf{Proof sketch.} Take any orthogonal $Q$ ($QQ'=I_n$, $Q\neq I$). Set $\Theta_\ell^* = \Theta_\ell Q$ and $\varepsilon_t^* = Q'\varepsilon_t$. Then:
    \[
    \sum_\ell \Theta_\ell \Theta_{\ell+h}' = \sum_\ell \Theta_\ell QQ' \Theta_{\ell+h}' = \sum_\ell (\Theta_\ell Q)(\Theta_{\ell+h} Q)' = \sum_\ell \Theta_\ell^* (\Theta_{\ell+h}^*)'.
    \]
    \textcolor{mitred}{$\Rightarrow$ We can always rotate the shocks and preserve the entire serial correlation structure.} 
\end{frame}

% ----------------------------------------------------------------------------
\begin{frame}{Shocks vs.\ Prediction Errors}
    \textbf{The Wold representation} (reducible-form):
    \[
    y_t = \mu + \sum_{\ell=0}^\infty B_\ell\, u_{t-\ell}, \qquad u_t \sim WN(0, \Sigma_u).
    \]
    \textbf{Under invertibility} (shocks are linear combinations of present and past data):
    \[
    u_t = H\varepsilon_t \implies \Sigma_u = HH', \quad B_\ell = \Theta_\ell H^{-1}.
    \]
    \begin{itemize}
        \item We can estimate $\{B_\ell\}$ and $\Sigma_u$ from data
        \item \textbf{But} this is not enough to recover $\{\Theta_\ell\}$
    \end{itemize}

    \vspace{0.4em}
    \textbf{The rotation parameterization.} Write $H = \text{Ch}(\Sigma_u)\cdot Q$ where $\text{Ch}(\Sigma_u)$ is the Cholesky square root:
    \[
    \Sigma_u = HH' \iff H = \text{Ch}(\Sigma_u)\,Q, \quad QQ' = I_n.
    \]
    \textcolor{mitred}{\textbf{Task:} impose restrictions on $\{\Theta_\ell\}$ to pin down a unique $Q$.}
\end{frame}

% ----------------------------------------------------------------------------
\begin{frame}{Identification Strategies}
    \textbf{Order condition.} $H$ has $n^2$ parameters but $\Sigma_u = HH'$ provides only $\tfrac{n(n+1)}{2}$ restrictions $\Rightarrow$ need $\tfrac{n(n-1)}{2}$ additional restrictions.

    \vspace{0.4em}
    \textbf{Two broad families:}
    \begin{enumerate}
        \item \textbf{Invertibility + restrictions} (classical VAR literature):
        \begin{itemize}
            \item Zero (short-run or long-run), sign, statistical
        \end{itemize}
        \item \textbf{Without invertibility} (recent):
        \begin{itemize}
            \item Proxy/Narrative SVARs, macro IVs
        \end{itemize}
    \end{enumerate}

    \vspace{0.4em}
    \textbf{This problem set covers:}
    \begin{itemize}
        \item \textbf{Problem 1}: Long-run zeros (Gal\'{\i} 1999) + short-run zeros (CEE 2005)
        \item \textbf{Problem 2}: Long-run restrictions in cointegrated systems (SVECM)
        \item \textbf{Problem 3}: \textcolor{mitred}{Critique} --- what goes wrong with Cholesky (CFP 2009)
    \end{itemize}
\end{frame}

% ============================================================================
% PROBLEM 1: GALÍ (1999)
% ============================================================================
\begin{frame}[plain]
    \begin{center}
        {\Large \textbf{Problem 1: Hybrid Identification}}\\[1em]
        {\large Gal\'{\i} (1999) Long-Run + CEE (2005) Short-Run Restriction}
    \end{center}
\end{frame}

% ----------------------------------------------------------------------------
\begin{frame}{Problem 1: Setup}
    \begin{block}{VAR and Structural Form}
        \vspace{-0.3em}
        \[
        y_t = (\Delta a_t,\; n_t,\; \pi_t,\; i_t)', \qquad y_t = \sum_{j=1}^p A_j y_{t-j} + u_t, \quad \E[u_t u_t'] = \Sigma_u.
        \]
        Structural form: $u_t = B_0^{-1} w_t$, $\E[w_t w_t'] = I_4$, $w_t = (\varepsilon_t^a, \varepsilon_t^d, \varepsilon_t^\pi, \varepsilon_t^m)'$.

        Structural MA: $y_t = \sum_{h=0}^\infty \Theta_h w_{t-h}$, $\;\Theta_0 = B_0^{-1}$.

        Long-run multiplier: $\Theta(1) = \sum_{h=0}^\infty \Theta_h = \Phi(1)\,B_0^{-1}$.
    \end{block}

    \vspace{0.3em}
    \textbf{Variables:}
    \begin{itemize}
        \item $\Delta a_t$: productivity growth \quad $n_t$: hours worked
        \item $\pi_t$: inflation \quad $i_t$: interest rate
        \item $\varepsilon_t^a$: technology \; $\varepsilon_t^d$: demand \; $\varepsilon_t^\pi$: cost-push \; $\varepsilon_t^m$: monetary
    \end{itemize}
\end{frame}

% ----------------------------------------------------------------------------
\begin{frame}{Problem 1(a): Gal\'{\i}'s Long-Run Restriction --- Derivation}
    \textbf{Gal\'{\i}'s assumption:} Only the technology shock $\varepsilon_t^a$ has a permanent effect on the \textit{level} of productivity $a_t$.

    \vspace{0.3em}
    \textbf{Mapping VAR growth rates to levels.} Since $a_T = a_{-1} + \sum_{t=0}^T \Delta a_t$, substitute the structural MA for $\Delta a_t$:
    \begin{align*}
        \frac{\partial a_T}{\partial \varepsilon_0^j}
        &= \sum_{t=0}^T \sum_{h=0}^\infty [\Theta_h]_{1,j}\,\mathbf{1}\{t - h = 0\}
        = \sum_{t=0}^T [\Theta_t]_{1,j}.
    \end{align*}
    Taking $T \to \infty$:
    \[
    \lim_{T\to\infty} \frac{\partial a_T}{\partial \varepsilon_0^j} = \sum_{t=0}^\infty [\Theta_t]_{1,j} = [\Theta(1)]_{1,j}.
    \]
    \textcolor{mitred}{The long-run effect on productivity \textit{levels} equals the $(1,j)$ entry of $\Theta(1)$.}
\end{frame}

% ----------------------------------------------------------------------------
\begin{frame}{Problem 1(a): Gal\'{\i}'s Long-Run Restriction --- Result}
    \begin{alertblock}{Identifying restriction}
        Only $\varepsilon_t^a$ ($j=1$) has a long-run effect on productivity levels. For $j \in \{2,3,4\}$:
        \[
        [\Theta(1)]_{1,2} = [\Theta(1)]_{1,3} = [\Theta(1)]_{1,4} = 0.
        \]
        Equivalently: $[\Theta(1)]_{1,\cdot} = \begin{pmatrix} * & 0 & 0 & 0 \end{pmatrix}$.
    \end{alertblock}

    \vspace{0.3em}
    \textbf{Connection to reduced form.} Since $\Theta(1) = \Phi(1)\,B_0^{-1} = A(1)^{-1} B_0^{-1}$:
    \[
    [\Theta(1)]_{1,j} = 0 \iff \text{first row of } A(1)^{-1}B_0^{-1} \text{ has zeros in cols 2,3,4.}
    \]

    \vspace{0.2em}
    \textbf{Three long-run restrictions} on the product $A(1)^{-1} B_0^{-1}$.
\end{frame}

% ----------------------------------------------------------------------------
\begin{frame}{Problem 1(b): CEE Short-Run Restriction}
    \textbf{Christiano--Eichenbaum--Evans (2005):} Productivity growth does not respond contemporaneously to a monetary policy shock:
    \[
    \frac{\partial(\Delta a_t)}{\partial \varepsilon_t^m}\bigg|_t = 0.
    \]
    The contemporaneous response equals the impact matrix entry:
    \[
    \frac{\partial(\Delta a_t)}{\partial \varepsilon_t^m}\bigg|_t = [B_0^{-1}]_{1,4} = 0.
    \]

    \begin{alertblock}{Short-run restriction}
        $[B_0^{-1}]_{1,4} = 0$: the $(1,4)$ entry of the impact matrix is zero.
    \end{alertblock}

    \textbf{Interpretation:} Within the period, a monetary shock cannot affect productivity growth (no contemporaneous supply-side channel).
\end{frame}

% ----------------------------------------------------------------------------
\begin{frame}{Problem 1(c): Identification Analysis}
    \textbf{Degrees of freedom.} $B_0^{-1}$ is $4\times 4$ (16 elements). Normalization $B_0^{-1}(B_0^{-1})' = \Sigma_u$ gives $\frac{4\times 5}{2} = 10$ restrictions. \textbf{Remaining d.o.f.:} $16 - 10 = 6$.

    \vspace{0.3em}
    \textbf{Restrictions available:}
    \begin{itemize}
        \item 3 long-run restrictions from (a): $[\Theta(1)]_{1,2}=[\Theta(1)]_{1,3}=[\Theta(1)]_{1,4}=0$
        \item 1 short-run restriction from (b): $[B_0^{-1}]_{1,4}=0$
        \item \textbf{Total: 4 restrictions $<$ 6 d.o.f.} $\Rightarrow$ \textcolor{mitred}{under-identified by the order condition}
    \end{itemize}

    \vspace{0.3em}
    \textbf{Rank condition.} Write $B(\theta) = \hat{L}Q(\theta)$ ($\hat{L}$ = Cholesky of $\hat\Sigma_u$, $Q(\theta)$ orthogonal). Stack restrictions as $f(\theta) = 0$:
    \[
    f(\theta) = \begin{pmatrix} (\hat{L}Q)_{1,4} \\ (\hat\Phi(1)\hat{L}Q)_{1,2} \\ (\hat\Phi(1)\hat{L}Q)_{1,3} \\ (\hat\Phi(1)\hat{L}Q)_{1,4} \end{pmatrix} = 0.
    \]
    Local identification requires $\rank\!\left(\partial f/\partial\theta'\right) = 4$ at the true $\theta_0$. \textbf{2 additional restrictions needed} for exact identification.
\end{frame}

% ----------------------------------------------------------------------------
\begin{frame}{Problem 1(d): Computational Procedure}
    \textbf{Core idea:} estimate reduced form by OLS, then search over rotations $Q(\theta)$ until $f(\theta) = 0$.

    \begin{enumerate}
        \item \textbf{Estimate VAR by OLS.} Obtain $(\hat{A}_1,\ldots,\hat{A}_p)$, residuals $\hat{u}_t$, and $\hat{\Sigma}_u$.
        \vspace{0.2em}
        \item \textbf{Compute reduced-form long-run multiplier.}
        \[
        \hat{A}(1) = I - \hat{A}_1 - \cdots - \hat{A}_p, \qquad \hat\Phi(1) = \hat{A}(1)^{-1}.
        \]
        \vspace{-0.4em}
        \item \textbf{Factor $\hat\Sigma_u$ and parameterize $Q$.} Compute Cholesky $\hat{L}$ ($\hat\Sigma_u = \hat{L}\hat{L}'$). Parameterize $Q(\theta)$ as Givens rotations 
        \vspace{0.2em}
        \item \textbf{Solve $f(\theta) = 0$.} With 2 additional restrictions: square $6\times 6$ system, use root-finder. Without: set-identified --- report range.
        \vspace{0.2em}
        \item \textbf{Recover structural objects.} $\hat{B}_0^{-1} = \hat{L}Q(\hat\theta)$, $\;\hat\Theta_h = \hat\Phi_h \hat{B}_0^{-1}$.
    \end{enumerate}

    % \vspace{0.2em}
    % \textcolor{mitred}{\textit{Optional:} Blanchard--Quah as initialization --- factorize $\hat\Phi(1)\hat\Sigma_u\hat\Phi(1)'$ to get a starting rotation.}
\end{frame}

% ----------------------------------------------------------------------------
\begin{frame}{Problem 1(e): Hours Fall After a Technology Shock}
    \textbf{Gal\'{\i}'s finding:} hours worked $n_t$ \textit{fall} after a positive technology shock --- contrary to standard RBC predictions.

    \vspace{0.3em}
    \textbf{Identified pattern:}
    \begin{itemize}
        \item $[\Theta_0]_{2,1} < 0$: hours fall on impact
        \item Hours may remain below steady state for several periods
        \item Long run: hours return to (or slightly above) steady state
        \item Inflation: $[\Theta_h]_{3,1} < 0$ for small $h$ --- technology reduces marginal cost
    \end{itemize}

    \vspace{0.3em}
    \textbf{Reconciliation with long-run restriction:}
    \begin{itemize}
        \item $[\Theta(1)]_{1,1} > 0$: technology permanently raises productivity (\checkmark)
        \item Mechanism: sticky prices $\Rightarrow$ firms cannot lower prices immediately $\Rightarrow$ reduce labor input
        \item In the long run, prices adjust, output expands, productivity remains permanently higher
    \end{itemize}
    \textcolor{mitred}{The long-run restriction on productivity does \textit{not} constrain the short-run response of hours.}
\end{frame}

% ============================================================================
% PROBLEM 2: SVECM
% ============================================================================
\begin{frame}[plain]
    \begin{center}
        {\Large \textbf{Problem 2: SVECM}}\\[1em]
        {\large Permanent and Transitory Shocks in a Cointegrated System}
    \end{center}
\end{frame}

% ----------------------------------------------------------------------------
% Background: Cointegration and VECMs (brief intro for students)
% ----------------------------------------------------------------------------
\begin{frame}{Background: Cointegration}
    \begin{block}{Definition}
        $x_t$ ($I(1)$, $n\times1$) is \textbf{cointegrated} if $\exists$ nonzero $\beta$ s.t.\ $\beta'x_t \sim I(0)$.
    \end{block}
    \pause
    \vspace{0.4em}
    \textbf{Why does this arise?} Macro variables grow (output, consumption, investment are all $I(1)$), but economic theory implies they share common trends --- e.g.\ balanced growth.
    \pause
    \vspace{0.4em}
    \textbf{Why do we care for structural analysis?}
    \begin{itemize}
        \item Some shocks are \textbf{permanently non-neutral} (technology); others are \textbf{neutral in the long run} (monetary)
        \item Cointegration rank $r$ $\Rightarrow$ exactly $n - r$ \textbf{common stochastic trends} $\Rightarrow$ at most $n-r$ permanent shocks
        \item This directly constrains what structural identification is possible
    \end{itemize}
    \pause
    \emph{We want to find a representation which is stationary and uses cointegration structure!}
\end{frame}

% ----------------------------------------------------------------------------
\begin{frame}{Background: How to Model $I(1)$ Cointegrated Data?}
    \vspace{-1em}
    What is the problem with the usual VAR?
    \pause 
    \begin{itemize}
        \item \textbf{VAR in levels} --- valid; consistent; does not exploit cointegration structure
        \item \textbf{VAR in first differences} --- $\Delta x_t$ is stationary, but:
    \end{itemize}
    \pause
    \begin{alertblock}{Fundamental problem (Beveridge--Nelson)}
       If $x_t\sim I(1)$, then $\Delta x_t$ has no VAR representation (despite stationary)
    \end{alertblock}
    \pause
    If we correct for non-invertibility, we get a VAR-like representation: the \textbf{VECM}.
    \vspace{0.3em}
    \begin{block}{VECM (follows from Granger Representation Theorem)}
    If $x_t\sim I(1)$ with cointegration rank $r$, then $\Delta x_t$ is stationary and can be represented as:
        \begin{equation*}
        \Delta x_t = \underbrace{\alpha\beta' x_{t-1}}_{\text{error correction}} + \sum_{j=1}^{p-1}\Gamma_j\,\Delta x_{t-j} + u_t
        \end{equation*}
        where $\beta$: cointegrating vectors; and $\alpha$: adjustment (loading) coefficients
    \end{block}
\end{frame}

% ----------------------------------------------------------------------------
\begin{frame}{Background: From VECM to SVECM}
    Add structural layer $u_t = B_0^{-1}w_t$, $\E[w_tw_t']=I_n$ --- exactly as in a standard SVAR.

    \vspace{0.4em}
    \textbf{What cointegration buys for identification:}
    \begin{itemize}
        \item Non-stationary long-run level representation:
        \begin{equation*}
        x_t = C\sum_{s=1}^t u_s + C^*(L)u_t + x_0^*,
        \end{equation*}
        where $C$ is the \textbf{reduced-form long-run multiplier} (cumulated impact on levels)
        \item \textbf{Structural long-run impact:} since $u_t=B_0^{-1}w_t$, the long-run impact of $w_t$ on $x_t$ is $\Upsilon = C\,B_0^{-1}$
        \item $\rank(C)=n-r$ $\Rightarrow$ only $n-r$ \textbf{permanent shocks}; transitory shocks have \textbf{zero columns in} $\Upsilon$
    \end{itemize}

    \vspace{0.3em}
    \textcolor{mitred}{This is the Shapiro--Watson (1988) / KPSW (1991) approach --- exactly what Problem 2 asks you to implement.}
\end{frame}

% ----------------------------------------------------------------------------
\begin{frame}{Problem 2: Setup}
    \vspace{-1em}
    \begin{block}{VECM for $x_t = (y_t, c_t, i_t)'$}
        \vspace{-0.5em}
        \[
        \Delta x_t = \alpha\beta' x_{t-1} + \sum_{j=1}^{p-1}\Gamma_j \Delta x_{t-j} + u_t, \qquad \E[u_t u_t'] = \Sigma_u.
        \]
        $x_t$ is $I(1)$, cointegration rank $r = 1$. $\;\alpha,\beta \in \R^{3\times 1}$.

        Structural shocks: $u_t = B_0^{-1}w_t$, $w_t = (\varepsilon_t^P, \varepsilon_t^T, \varepsilon_t^M)'$, $\E[w_tw_t'] = I_3$.
    \end{block}
    \vspace{-0.5em}
    The VECM implies (\textbf{Granger Representation Theorem}):
    \[
    x_t = C\sum_{s=1}^t u_s + C^*(L)u_t + x_0^*, \qquad C = \beta_\perp\bigl(\alpha_\perp'\Gamma(1)\beta_\perp\bigr)^{-1}\alpha_\perp'.
    \]
    \begin{itemize}
        \item Orthogonal Complement: $\beta^\prime \beta_\perp = 0$
        \item $C\sum u_s$: stochastic trends (permanent component)
        \item $C^*(L)u_t$: transitory dynamics
        \item $\rank(C) = n - r = 2$: exactly 2 common stochastic trends
    \end{itemize}
\end{frame}

% ----------------------------------------------------------------------------
\begin{frame}{Problem 2(a): Common Stochastic Trends}
    \begin{alertblock}{Result}
        With $n = 3$ variables and cointegration rank $r = 1$:
        \[
        \text{Number of common stochastic trends} = n - r = 3 - 1 = 2.
        \]
    \end{alertblock}

    \vspace{0.3em}
    \textbf{Implication for structural identification:}
    \begin{itemize}
        \item Since $\rank(C) = 2$, at most \textbf{2 shocks can have permanent effects} on the levels
        \item The third shock must be \textbf{transitory} (no long-run effect on any variable)
        \item This is \textcolor{mitred}{not an identifying assumption} --- it is a consequence of cointegration
        \item It provides 1 free restriction; we need 1 more for exact identification
    \end{itemize}
\end{frame}

% ----------------------------------------------------------------------------
\begin{frame}{Problem 2(b)--(c): Long-Run Structural Restrictions}
    \textbf{Structural long-run impact matrix.} Substituting $u_t = B_0^{-1}w_t$ into the Granger representation:
    \[
    x_t = \underbrace{CB_0^{-1}\sum_{s=1}^t w_s}_{\text{permanent}} + \underbrace{C^*(L)B_0^{-1}w_t}_{\text{transitory}} + x_0^*.
    \]
    \[
    \Upsilon \equiv CB_0^{-1}: \quad \text{structural long-run impact matrix on levels.}
    \]

    \begin{alertblock}{(b) Transitory shock $\varepsilon_t^T$ (Shapiro--Watson / KPSW)}
        $[\Upsilon]_{\cdot,2} = \mathbf{0}$: no long-run effect on output, consumption, or investment.
        \textit{Note:} $\rank(C)=2$ makes only 2 of these 3 restrictions independent.
    \end{alertblock}

    \begin{alertblock}{(c) Monetary neutrality (long run)}
        $[\Upsilon]_{1,3} = 0$: monetary/financial shock has no long-run effect on output.
    \end{alertblock}
\end{frame}

% ----------------------------------------------------------------------------
\begin{frame}{Problem 2(d): Short-Run Restriction and Identification}
    \begin{alertblock}{(d) Short-run restriction (CEE-style)}
        Monetary shock does not affect output growth contemporaneously:
        \[
        [B_0^{-1}]_{1,3} = 0.
        \]
    \end{alertblock}

    \vspace{0.3em}
    \textbf{Identification count:}
    \begin{itemize}
        \item $B_0^{-1}$ is $3\times 3$: $9 - 6 = 3$ d.o.f.
        \item Restrictions: 2 independent (from (b)) + 1 (from (c)) + 1 (from (d)) = \textbf{4}
        \item \textbf{Over-identified by the count} --- need to verify consistency
    \end{itemize}

    \vspace{0.3em}
    \textbf{Key feasibility check.} Since $\rank(C)=2$, column 2 of $\Upsilon$ is zero iff $[B_0^{-1}]_{\cdot,2}$ lies in $\ker(C)$ (1-dimensional). This restricts column 2 of $B_0^{-1}$ to a line --- numerical verification required for specific $(\hat{C},\hat\Sigma_u)$.
\end{frame}

% ----------------------------------------------------------------------------
\begin{frame}{Problem 2(f): Why a VAR in First Differences Is Invalid}
    \textbf{Wold representation of $\Delta x_t$:}
    \[
    \Delta x_t = \delta + \psi(L)\varepsilon_t, \qquad \psi(1) = \text{long-run multiplier matrix.}
    \]

    \begin{alertblock}{Key result (Beveridge--Nelson / Granger)}
        Cointegration ($r=1$) $\Rightarrow$ $\rank(\psi(1)) = n - r = 2 < 3$ $\Rightarrow$ $|\psi(1)|=0$.

        \medskip
        $L=1$ is a root of $|\psi(L)|=0$ $\Rightarrow$ unit MA root $\Rightarrow$ $\psi(L)$ is \textbf{noninvertible}.

        \medskip
        \textbf{$\Delta x_t$ does not have a VAR representation.}
    \end{alertblock}

    \vspace{0.3em}
    \textcolor{mitred}{A finite-order VAR in first differences is fundamentally invalid for a cointegrated system---one cannot ``flip'' a unit root on the unit circle.}
\end{frame}

% ----------------------------------------------------------------------------
\begin{frame}{Problem 2(f): SVECM vs.\ VAR in Levels}
    \vspace{-1em}
    \textbf{The correct comparison:} SVECM vs.\ unrestricted VAR in levels ($x_t = \sum_j A_j x_{t-j} + u_t$).

    \vspace{0.3em}
    \begin{block}{SVECM}
        \begin{itemize}
            \item[\textcolor{darkblue}{+}] Exploits reduced-rank structure $\alpha\beta'$ $\Rightarrow$ more efficient
            \item[\textcolor{mitred}{--}] \textbf{Requires correct specification of $r$}
        \end{itemize}
    \end{block}

    \begin{block}{VAR in levels}
        \begin{itemize}
            \item[\textcolor{darkblue}{+}] Consistent regardless of true $r$ (Stock, 1987); no rank commitment needed
            \item[\textcolor{mitred}{--}] Less efficient; non-standard inference on individual coefficients
        \end{itemize}
    \end{block}

    \vspace{0.3em}
    \textbf{Failure modes (SVECM):}
    \begin{itemize}
        \item \textbf{$r$ overspecified} ($r=2$, true $r=1$): $\rank(C)$ too low $\Rightarrow$ $\Upsilon$ inconsistent
        \item \textbf{$r$ underspecified} ($r=0$): reduces to VAR in $\Delta x_t$ $\Rightarrow$ invalid (noninvertible MA)
    \end{itemize}
\end{frame}

% ============================================================================
% PROBLEM 3: CFP CRITIQUE
% ============================================================================
\begin{frame}[plain]
    \begin{center}
        {\Large \textbf{Problem 3: The CFP Critique}}\\[1em]
        {\large Recursive Identification Under Misspecification}
    \end{center}
\end{frame}

% ----------------------------------------------------------------------------
\begin{frame}{Problem 3: Setup}
    \begin{block}{True DGP and Econometrician's Model}
        True DGP: $z_t = \sum_{h=0}^\infty \Psi_h \varepsilon_{t-h}$, $\;\varepsilon_t = (\varepsilon_t^a, \varepsilon_t^c, \varepsilon_t^m)'$, $\;\E[\varepsilon_t\varepsilon_t'] = I_3$.

        Econometrician: estimates VAR, imposes Cholesky ordering $(y,\pi,i)$.

        $u_t = B_0^{-1} w_t$ with $B_0^{-1}$ \textbf{lower triangular}.
    \end{block}

    \vspace{0.3em}
    \textbf{Key assumption:} In the true DGP, monetary policy affects output \textit{contemporaneously}:
    \[
    (\Psi_0)_{1,3} \neq 0 \qquad \text{(true DGP)}.
    \]
    But Cholesky imposes:
    \[
    (B_0^{-1})_{1,3} = 0 \qquad \text{(identification assumption)}.
    \]
    \textcolor{mitred}{These two conditions are contradictory --- this is the source of misidentification.}
\end{frame}

% ----------------------------------------------------------------------------
\begin{frame}{Problem 3(a): Cholesky Zero Restrictions}
    With ordering $(y, \pi, i)$, the lower-triangular $B_0^{-1}$ imposes:
    \[
    B_0^{-1} = \begin{pmatrix} b_{11} & 0 & 0 \\ b_{21} & b_{22} & 0 \\ b_{31} & b_{32} & b_{33} \end{pmatrix}.
    \]

    \vspace{0.3em}
    \textbf{Which variables respond to which shocks (contemporaneously):}
    \begin{center}
    \begin{tabular}{l|ccc}
    \toprule
    & $w_t^1$ & $w_t^2$ & $w_t^3$ (``monetary'') \\
    \midrule
    $y_t$ & \checkmark & $\times$ & $\times$ \\
    $\pi_t$ & \checkmark & \checkmark & $\times$ \\
    $i_t$ & \checkmark & \checkmark & \checkmark \\
    \bottomrule
    \end{tabular}
    \end{center}

    \vspace{0.2em}
    \textbf{Economic content:} the monetary shock affects only the interest rate on impact --- neither output nor inflation respond within the period.
\end{frame}

% ----------------------------------------------------------------------------
\begin{frame}{Problem 3(b): Non-Uniqueness and Cholesky}
    \textbf{The non-uniqueness problem.} $\Sigma_u = B_0^{-1}(B_0^{-1})'$ has $\tfrac{3\cdot 4}{2}=6$ equations in 9 unknowns $\Rightarrow$ 3 d.o.f.\ remain. For any orthogonal $Q$:
    \[
    \tilde{B}_0^{-1} = B_0^{-1}Q \implies \tilde{B}_0^{-1}(\tilde{B}_0^{-1})' = B_0^{-1}QQ'(B_0^{-1})' = \Sigma_u.
    \]
    Infinitely many valid impact matrices.

    \vspace{0.3em}
    \textbf{How Cholesky selects a rotation.} Let $\Sigma_u^{1/2}$ be the lower-triangular Cholesky square root. Cholesky sets:
    \[
    B_0^{-1} = \Sigma_u^{1/2} \iff Q = I \quad \text{in the representation } B_0^{-1} = \Sigma_u^{1/2}Q.
    \]
    \begin{itemize}
        \item This is an \textbf{economic assumption about timing}, not a statistical necessity
        \item Different variable orderings yield different identified shocks
        \item \textcolor{mitred}{Assuming $Q=I$ is valid \textit{only if} $B_0^{-1} = \Psi_0$ (correct identification)}
    \end{itemize}
\end{frame}

% ----------------------------------------------------------------------------
\begin{frame}{Problem 3(c): Identified Shock as a Mixture --- Derivation}
    \textbf{Step 1: Reduced-form innovation = $\Psi_0\varepsilon_t$.} From the structural MA $z_t = \sum_{h\geq0}\Psi_h\varepsilon_{t-h}$, under invertibility:
    \[
    u_t = z_t - \mathrm{Proj}(z_t \mid \mathcal{F}^z_{t-1}) = \sum_{h\geq0}\Psi_h\varepsilon_{t-h} - \sum_{h\geq1}\Psi_h\varepsilon_{t-h} = \Psi_0\varepsilon_t.
    \]

    \textbf{Step 2: Econometrician applies Cholesky.} $w_t = B_0 u_t = B_0\Psi_0\varepsilon_t$. Define:
    \[
    R \equiv B_0\Psi_0, \qquad w_t = R\varepsilon_t.
    \]
    Under correct identification: $B_0^{-1} = \Psi_0 \Rightarrow R = I$. \textcolor{mitred}{Here $R \neq I$.}

    \vspace{0.3em}
    \textbf{Step 3: Proof that $w_t^m$ is a mixture.} If $R_{31}=R_{32}=0$, then $[B_0^{-1}]_{\cdot,3} = [\Psi_0]_{\cdot,3}$. But $[B_0^{-1}]_{1,3}=0$ (Cholesky) and $[\Psi_0]_{1,3}\neq0$ (true DGP) --- \textbf{contradiction}.

    \begin{alertblock}{Result}
        $w_t^m = R_{31}\varepsilon_t^a + R_{32}\varepsilon_t^c + R_{33}\varepsilon_t^m$, with $(R_{31},R_{32})\neq(0,0)$.
    \end{alertblock}
\end{frame}

% ----------------------------------------------------------------------------
\begin{frame}{Problem 3(d): Sufficient Condition for Sign Reversal --- Step 1}
    \textbf{IRF linearity.} By linearity of impulse responses:
    \[
    \text{IRF}_y^{w^m}(h) = R_{31}\cdot\text{IRF}_y^{\varepsilon^a}(h) + R_{32}\cdot\text{IRF}_y^{\varepsilon^c}(h) + R_{33}\cdot\text{IRF}_y^{\varepsilon^m}(h).
    \]

    \textbf{Evaluate at $h=0$.} Since $\text{IRF}_y^{\varepsilon^j}(0) = (\Psi_0)_{1,j}$:
    \[
    \text{IRF}_y^{w^m}(0) = R_{31}(\Psi_0)_{1,1} + R_{32}(\Psi_0)_{1,2} + R_{33}(\Psi_0)_{1,3}.
    \]
    But Cholesky forces $(B_0^{-1})_{1,3} = 0$, so $\text{IRF}_y^{w^m}(0) = 0$. Therefore:
    \[
    R_{33}\,(\Psi_0)_{1,3} = -\bigl[R_{31}(\Psi_0)_{1,1} + R_{32}(\Psi_0)_{1,2}\bigr]. \tag{$*$}
    \]

    \begin{block}{Remarks}
        \begin{itemize}
            \item $(^*)$ is \textbf{automatic} --- it holds for any $\Psi_0$, not an extra assumption
            \item It shows: at $h=0$, the monetary term and contamination term are \textit{equal and opposite}
            \item Sign reversal at $h=0$ is \textbf{impossible by construction}; it manifests at $h\geq1$
        \end{itemize}
    \end{block}
\end{frame}

% ----------------------------------------------------------------------------
\begin{frame}{Problem 3(d): Sufficient Condition for Sign Reversal --- Step 2}
    \begin{alertblock}{Sufficient condition (two parts)}
        \textbf{1. Restriction on $\Psi_0$ (contemporaneous impact matrix):}
        \[
        (\Psi_0)_{1,3} < 0 \quad \text{(true monetary shock lowers output on impact).}
        \]
        Since $R_{33}>0$, constraint $(^*)$ implies $R_{31}(\Psi_0)_{1,1}+R_{32}(\Psi_0)_{1,2}>0$: the contamination enters with the \textit{opposite} sign, hidden at $h=0$ by exact cancellation.

        \vspace{0.4em}
        \textbf{2. Dynamic restriction:} $\exists\, h\geq1$ such that $(\Psi_h)_{1,3}<0$ and:
        \[
        R_{31}(\Psi_h)_{1,1} + R_{32}(\Psi_h)_{1,2} > |R_{33}(\Psi_h)_{1,3}|.
        \]
    \end{alertblock}

    \vspace{0.2em}
    \textbf{Conclusion.} Under these conditions, $\text{IRF}_y^{w^m}(h) > 0$ while $\text{IRF}_y^{\varepsilon^m}(h) < 0$: \textcolor{mitred}{the identified IRF has the wrong sign.}

    % Note: $R_{3j}$ are determined by $\Psi_0$ through $\Sigma_u = \Psi_0\Psi_0'$ --- both conditions are ultimately restrictions on the true DGP.
\end{frame}

% ----------------------------------------------------------------------------
\begin{frame}{Problem 3(e): Two Robustness Strategies}
    \textbf{Strategy 1: Sign Restrictions} \textcolor{mitred}{(Uhlig 2005)}
    \begin{itemize}
        \item Replace exact zeros with \textit{inequality} constraints on IRFs:
        \begin{itemize}
            \item $\text{IRF}_i^{\varepsilon^m}(0) > 0$, $\;\text{IRF}_y^{\varepsilon^m}(h) < 0$, $\;\text{IRF}_\pi^{\varepsilon^m}(h) < 0$
        \end{itemize}
        \item \textbf{Identifies}: a \textit{set} of valid rotations $\Rightarrow$ set-identified IRFs
        \item \textbf{Advantage}: robust to timing misspecification
    \end{itemize}

    \vspace{0.4em}
    \textbf{Strategy 2: Proxy SVAR / External Instruments} \textcolor{mitred}{(Gertler--Karadi 2015, Romer--Romer 2004)}
    \begin{itemize}
        \item Use instrument $z_t$ correlated with $\varepsilon_t^m$ and uncorrelated with $(\varepsilon_t^a, \varepsilon_t^c)$
        \item IRF identified (up to scale) via IV:
        \[
        \text{IRF}_y^{\varepsilon^m}(h) \propto \frac{\Cov(y_{t+h}, z_t)}{\Cov(i_t, z_t)}.
        \]
        \item \textbf{Advantage}: no timing assumptions; identification from external information
    \end{itemize}
\end{frame}

% ============================================================================
% SUMMARY
% ============================================================================
\begin{frame}{Summary}
    % \begin{center}
    % \begin{tabular}{l|l|l}
    % \toprule
    % \textbf{Problem} & \textbf{Key restriction} & \textbf{Identifies} \\
    % \midrule
    % 1: Gal\'{\i} (1999) & $[\Theta(1)]_{1,j}=0$, $j\neq1$ + $[B_0^{-1}]_{1,4}=0$ & Technology shock IRFs \\
    % \midrule
    % 2: SVECM & $[\Upsilon]_{\cdot,2}=0$, $[\Upsilon]_{1,3}=0$, $[B_0^{-1}]_{1,3}=0$ & Permanent vs.\ transitory \\
    % \midrule
    % 3: CFP & \textit{Critique} of Cholesky & Sign restrictions, Proxy SVAR \\
    % \bottomrule
    % \end{tabular}
    % \end{center}

    % \vspace{0.5em}
    \textbf{Big picture:}
    \begin{itemize}
        \item All identification strategies are about pinning down the rotation $Q$ in $H = \text{Ch}(\Sigma_u)Q$
        \item Short-run and long-run restrictions are two complementary tools
        \item Misspecified restrictions (like Cholesky when $(\Psi_0)_{1,3}\neq0$) can flip the sign of IRFs
        \item \textcolor{mitred}{Always check economic plausibility of identifying assumptions}
    \end{itemize}
\end{frame}

\end{document}
